%% P12 --- Combinatorics and counting
%%
%% Every number in this program that the reader cannot do in their head comes
%% from code/p12_combinatorics.py and is pulled in with \val{}. The console
%% listing is a committed transcript written by that script.
%%
%% THE THESIS: two questions -- does order matter, may anything repeat --
%% decide which of four counts applies, and almost every counting mistake an
%% engineer makes is answering the wrong one of the two.
%%
%% WHAT THIS PROGRAM IS OWED, read out of the files rather than remembered:
%%   F10  is the elementary layer under this one and says so in its own text.
%%        It gives the product rule with its independence condition, 2^n
%%        subsets, the pair count n(n-1)/2 with the "halve it" step, and
%%        inclusion-exclusion for two sets -- and defers "the four rules
%%        formally, pigeonhole, the birthday calculation" to here by name.
%%        THREE of its committed values are gated against in the script.
%%   F04  gives the factorial as a pi with the index in the body, and the
%%        empty product, which is why 0! = 1 needs no separate rule.
%%   F03  gives logarithms, the only way to hold 32000^20 in a head.
%%   P02  section 4 is "Never form a probability", and section 4 here is that
%%        finding arriving where nobody expects it: the textbook collision
%%        formula returns EXACTLY ZERO in float64 at 128 bits.
%%
%% WHAT THIS PROGRAM LEAVES ALONE, checked against tools/programs.json:
%%   what a probability IS, and conditioning                        -> P23
%%   random variables, expectation, variance                        -> P24
%%   graphs and their counting arguments                            -> P13
%%   quantifiers, induction, writing a proof                        -> P14
%%   O-notation. This program counts; it never says "quadratic"     -> P03
%%
%% Twin: programs/pl/P12-combinatorics.tex. Frame for frame, macro for macro,
%% number for number; tools/parity.py compares them.

\program{Combinatorics and counting}
\label{prog:P12}
\index{counting!four rules}
\index{combinatorics}

\begin{outcomes}
\outcome{1--8}{ask the two questions that decide which counting rule applies,
  and say why a count with neither answered is not a question yet}
\outcome{9--16}{compute arrangements and choices, and say where the factorial
  in the denominator of a binomial coefficient comes from}
\outcome{17--26}{use inclusion--exclusion beyond two sets, and say what the
  pigeonhole principle guarantees that no probability can}
\outcome{27--35}{size a hash against the number of items it must keep apart,
  and say which of two formulas for that is safe to evaluate}
\outcome{36--43}{put a number on a search space and on an exact attribution,
  and say what that number rules out}
\end{outcomes}

\begin{quiz}
\item From $\val{p12.four.n}$ features you pick $\val{p12.four.k}$ to ablate.
  How many ways? \teachesat{1--2}
  \answerto{The question is not complete. Until it says whether order matters
  and whether a feature may be picked twice it has three defensible answers,
  and the whole of section 1 is that pair of questions.}
\item What is $\binom{\val{p12.sym.n}}{\val{p12.sym.k}}$ compared with
  $\binom{\val{p12.sym.n}}{\val{p12.sym.rest}}$? \teachesat{13--14}
  \answerto{Equal, both $\val{p12.sym.val}$. Choosing $\val{p12.sym.k}$ to
  keep is choosing $\val{p12.sym.rest}$ to drop, so the two are the same act
  described from opposite ends.}
\item Three evaluation sets overlap pairwise and all three share some
  examples. Write the size of their union. \teachesat{21--22}
  \answerto{Add the three, subtract the three pairwise overlaps, add the
  triple overlap back. The last term is there because subtracting the pairs
  removed the triple region once too often.}
\item A hash has $\val{p12.bits64}$ bits and you have
  $\val{p12.hash.m}$ documents. Roughly what is the chance two collide?
  \teachesat{32--33}
  \answerto{About $\val{p12.hash.p64}$ per cent, which is
  $\val{p12.hash.ratio}$ times the figure most people reason their way to.
  The quantity that matters is the number of \emph{pairs}, not the number of
  documents.}
\item Exact Shapley values for $\val{p12.shap.big}$ features. How many
  coalitions must be evaluated? \teachesatone{40}
  \answerto{$\val{p12.shap.big.coal}$, which is $2^{\val{p12.shap.big}}$ and
  not $\val{p12.shap.big}!$. Only \emph{which} features came before matters,
  never their order, and collecting the orderings is what turns one
  exponential into a smaller one.}
\end{quiz}

% ============================================================
\section{Two questions before any formula}
\label{sec:P12-two-questions}
\index{counting!order and repetition}

\begin{fr}
Part~IV is a change of subject rather than a continuation. Nothing from
Program~\ref{prog:P04} to Program~\ref{prog:P11} is needed here, and nothing
here needs a matrix.

What this program is owed is Program~\ref{prog:F10}, which counted three
things \dash{} runs in a grid, subsets of a set, pairs from a collection
\dash{} and then named, in as many words, what it was leaving behind.

Start where it stopped. You have $\val{p12.four.n}$ features and you are going
to ablate $\val{p12.four.k}$ of them. How many ways are there to do that?
\dotline
\nextframe
\end{fr}

\begin{fr}
\begin{ansblock}
The question is not finished, and if you produced a number it was one of
three.
\end{ansblock}

Before any formula, two questions:

\begin{itemize}
\item \textbf{Does the order matter?} Ablating features $3$, $7$ and $1$ in
  that order, and ablating them in any other order, may or may not be the same
  experiment. The sentence did not say.
\item \textbf{May anything repeat?} A feature cannot be ablated twice in one
  run. A token can certainly be drawn twice from a vocabulary. The sentence
  did not say that either.
\end{itemize}

\textbf{Almost every counting mistake in engineering is answering the wrong
one of those two}, and the reason they are so easy to skip is that in ordinary
speech \enquote{pick three of ten} sounds like a complete instruction.

Take them one at a time. Suppose order does matter and repeats are allowed
\dash{} you are drawing $\val{p12.four.k}$ times from
$\val{p12.four.n}$ and writing down what you get. How many outcomes?
\nextframe
\end{fr}

\begin{fr}
\ans{$\val{p12.four.ordrep}$}
Which is $\val{p12.four.n}^{\val{p12.four.k}}$, and it is
Program~\ref{prog:F10}'s product rule with nothing added: three independent
choices of $\val{p12.four.n}$ each.

In general \textbf{$n^{k}$}, and it is the count you want whenever the thing
being counted is a \emph{sequence drawn from a fixed alphabet} \dash{} a
password, a $k$-token prompt, a byte string.

Now forbid the repeats. Order still matters, but once a feature is used it is
gone. How many?
\dotline
\nextframe
\end{fr}

\begin{fr}
\ans{$\val{p12.four.ordnorep}$}
From $\val{p12.four.n} \times 9 \times 8$: ten for the first, nine for the
second because one has gone, eight for the third.

That number has a name \dash{} it is the count of \textbf{arrangements}, or
permutations, of $\val{p12.four.k}$ things chosen from $\val{p12.four.n}$
\dash{} and it is worth a moment, because it looks as though it breaks the
rule that produced it.

Program~\ref{prog:F10} said independent choices multiply, and was careful to
say what independent meant: every choice of the first leaves the same options
open for the second. Here it does not. Pick feature $3$ first and feature $3$
is unavailable second. So does the product rule apply or not?
\nextframe
\end{fr}

\begin{fr}
\begin{ansblock}
It applies, and the condition it really needs is weaker than the one that was
stated.
\end{ansblock}

What has to be fixed is \textbf{the number of options at each step, not the
options themselves}. Whichever feature went first, exactly nine remain, so
every one of the ten first choices branches into the same nine. The tree is
uniform even though its branches carry different labels, and a uniform tree is
all the product rule ever needed.

\begin{note}
This is a refinement of a rule you already have rather than a new rule.
Program~\ref{prog:F10}'s statement is the one to use when checking a grid,
because there the options really are the same each time and a forbidden
combination breaks the count. This one is what lets a shrinking pool multiply
anyway, and the two are worth holding apart: the first is about which options
exist, the second only about how many.
\end{note}

Now drop the order. Ablating $\{3, 7, 1\}$ is one experiment however you list
it. How many now, and what did you divide by?
\dotline
\nextframe
\end{fr}

\begin{fr}
\ans{$\val{p12.four.unordnorep}$, dividing by $\val{p12.kfact}$}
Because each selection of $\val{p12.four.k}$ features was counted once for
every order it could have been written in, and $\val{p12.four.k}$ things can
be written in $\val{p12.four.k}! = \val{p12.kfact}$ orders.

\textbf{That is Program~\ref{prog:F10}'s \enquote{halve it} generalised.} That
program counted pairs by multiplying $n$ by $n-1$ and then halving, saying the
halving was a set fact: $\{a, b\}$ and $\{b, a\}$ are one pair. The halving
was a division by $2!$. With three at a time it is a division by $3!$, and
with $k$ at a time by $k!$ \dash{} the same step, and the only reason it did
not look like one is that dividing by two rarely looks like dividing by a
factorial.
\end{fr}

\begin{fr}
That leaves one cell of the table empty: order does not matter, and repeats
\emph{are} allowed. Counting bags rather than sequences \dash{} how many
different multisets of $\val{p12.four.k}$ tokens can be drawn from a
vocabulary of $\val{p12.four.n}$, if you keep no record of the order.

The answer is $\binom{n + k - 1}{k}$, which here is $\val{p12.four.unordrep}$.

\begin{rigourbox}
That formula is stated and not derived. The argument is the one called
\emph{stars and bars} and it is in any first combinatorics course: line up
$k$ items and $n-1$ dividers, and every arrangement of the two kinds of symbol
is one multiset.

It is the least used of the four in practice, which is why it is here as a
name to recognise rather than as a technique. When it does come up it is
usually disguised: counting the ways a fixed budget can be split between
$n$ buckets is this cell, and so is counting the shapes of a bag-of-words
vector with a fixed total count.
\end{rigourbox}
\end{fr}

\begin{fr}
The whole of the section, then, is one table.

\begin{center}
\begin{tabular}{lll}
\toprule
 & \textbf{repeats allowed} & \textbf{no repeats} \\
\midrule
\textbf{order matters} & $n^{k}$ \dash{} $\val{p12.four.ordrep}$
  & $\frac{n!}{(n-k)!}$ \dash{} $\val{p12.four.ordnorep}$ \\
\textbf{order does not} & $\binom{n+k-1}{k}$ \dash{} $\val{p12.four.unordrep}$
  & $\binom{n}{k}$ \dash{} $\val{p12.four.unordnorep}$ \\
\bottomrule
\end{tabular}
\end{center}

Four counts of the same ten things taken three at a time, and they run from
$\val{p12.four.unordnorep}$ to $\val{p12.four.ordrep}$ \dash{} a factor of
more than eight, decided entirely by two words that were never in the
question.

\mermaidfig{p12-two-questions}{The two questions that pick out which of four
counts a problem wants.}{order, repeats, which count}
\end{fr}

% ============================================================
\section{Arrangements and choices}
\label{sec:P12-choices}
\index{combinatorics!binomial coefficient}

\begin{fr}
Take the top-right cell to its limit. Instead of arranging
$\val{p12.four.k}$ of the $\val{p12.four.n}$ features, arrange all of them.

You have the tool for this already: Program~\ref{prog:F04} wrote the factorial
as a pi with the index in its own body, $n! = \prod_{i=1}^{n} i$. How many
orderings of $\val{p12.four.n}$ features are there?
\nextframe
\end{fr}

\begin{fr}
\ans{$\val{p12.fact.n}$}
Which is $\val{p12.four.n}!$, and it is already a number nobody would
enumerate.

Notice how it relates to the arrangement count. An ordering of all ten is an
arrangement of the first three followed by an ordering of the remaining seven,
so
\[ n! = \frac{n!}{(n-k)!} \times (n-k)! \]
which is the shrinking pool read backwards, and is why the arrangement count
is usually written as that quotient rather than as a product of $k$ shrinking
factors. The quotient is the tidier form; the product is the one you can do in
your head.
\end{fr}

\begin{fr}
The bottom-right cell has a symbol and a name, and both are worth having
because they are what the literature uses.

\[ \binom{n}{k} = \frac{n!}{k!\,(n-k)!} \]

the \textbf{binomial coefficient}: arrangements of $k$ from $n$, divided by
the $k!$ orderings each selection was counted in.

\begin{notationbox}
In English this is read \enquote{$n$ choose $k$}. In Polish it is
\emph{symbol Newtona} and is read \emph{$n$ po $k$} \dash{} a different name
for the same object, not a different object, and the Polish edition of this
book uses the Polish name throughout.

The notation itself is the same in both, which is the general rule this book
follows: what you would type stays the same, what you would say may not.
\end{notationbox}

Now use it on something you have already computed by another route.
Program~\ref{prog:F10} counted the pairs from $\val{p12.ie.a}$ documents. Write
that count as a binomial coefficient.
\dotline
\nextframe
\end{fr}

\begin{fr}
\ans{$\binom{\val{p12.ie.a}}{2} = \val{f10.pairs}$}
A pair is two things chosen from many, with the order discarded, which is the
bottom-right cell with $k = 2$. Expand it and
$\frac{n!}{2!\,(n-2)!} = \frac{n(n-1)}{2}$, which is the formula that program
gave.

\begin{note}
This is not a resemblance. \code{code/p12\_combinatorics.py} reads
Program~\ref{prog:F10}'s committed pair count out of its values file and fails
the build if $\binom{n}{2}$ does not reproduce it, so the two programs cannot
drift into disagreeing about what a pair is. Two more of that program's
numbers are gated the same way in the sections below.
\end{note}
\end{fr}

\begin{fr}
Here is the property of the binomial coefficient that people find hardest to
believe, and it is worth being wrong about first.

From $\val{p12.sym.n}$ features, is it easier to choose $\val{p12.sym.k}$ of
them or to choose $\val{p12.sym.rest}$ of them? Answer quickly, with no
working.
\nextframe
\end{fr}

\begin{fr}
\begin{ansblock}
Neither. Both are $\val{p12.sym.val}$.
\end{ansblock}

\begin{trapbox}
\textbf{$\binom{n}{k} = \binom{n}{n-k}$, and the reason is one sentence:
choosing which $\val{p12.sym.k}$ to keep \emph{is} choosing which
$\val{p12.sym.rest}$ to drop.} There is one act, and the two counts describe
it from opposite ends.

The reason almost everyone says the smaller number is easier is that the two
questions feel like different amounts of work \dash{} naming three things is
quicker than naming seventeen. That is a fact about writing them down, not
about how many there are. The formula makes it obvious, since swapping $k$ and
$n-k$ swaps two factors in a denominator that is a product.

Where this stops being a curiosity is at the ends. $\binom{n}{0}$ and
$\binom{n}{n}$ are both $1$ \dash{} there is exactly one way to choose nothing
and exactly one way to choose everything \dash{} and the first of those is the
empty product Program~\ref{prog:F04} argued for, arriving as a count.
\end{trapbox}
\end{fr}

\begin{fr}
So the counts rise from the ends and peak in the middle:
$\binom{\val{p12.sym.n}}{10} = \val{p12.sym.peak}$, the largest entry in that
row.

That shape decides a real experiment. You have $\val{p12.abl.n}$ features and
you want an ablation study. Testing \emph{every} subset is
Program~\ref{prog:F10}'s $2^{n}$, which is $\val{p12.abl.all}$ runs and is not
happening.

So restrict it: no subset larger than $\val{p12.abl.k}$. How many runs is
that?
\dotline
\nextframe
\end{fr}

\begin{fr}
\ans{$\val{p12.abl.upto}$}
From $1 + \val{p12.abl.singles} + \val{p12.abl.pairs} + \val{p12.sym.val}$
\dash{} the empty ablation, the singles, the pairs and the triples, which are
$\binom{20}{0}$ through $\binom{20}{3}$ and include the $\val{p12.sym.val}$
from two frames ago.

\textbf{$\val{p12.abl.upto}$ runs against $\val{p12.abl.all}$ is a factor of
$\val{p12.abl.frac}$}, and that is the entire content of the decision to cap
the subset size. It is also why the cap is usually two or three rather than
five: the row is still climbing steeply there, so each extra size costs more
than all the sizes below it put together.

The habit is Program~\ref{prog:F10}'s, one rule further on: write the count
down before writing the loop, and now you can write it down for any of the
four shapes rather than only for subsets.
\end{fr}

% ============================================================
\section{The rules that are not products}
\label{sec:P12-not-products}
\index{counting!inclusion--exclusion}
\index{counting!pigeonhole principle}

\begin{fr}
Everything so far has been the product rule wearing different clothes. Three
things are not, and each is the answer to a question a product cannot reach.

Start with the cheapest. Fix one particular feature \dash{} call it the first
one \dash{} and split every $k$-subset of $n$ into two piles by whether that
feature is in it. Count the two piles.
\nextframe
\end{fr}

\begin{fr}
\begin{ansblock}
$\binom{n-1}{k-1}$ contain it and $\binom{n-1}{k}$ do not, so
\[ \binom{n}{k} = \binom{n-1}{k-1} + \binom{n-1}{k} \]
\end{ansblock}

If the feature is in, the other $k-1$ come from the remaining $n-1$; if it is
out, all $k$ do. Nothing is in both piles and nothing is in neither, so the
counts add.

\textbf{That is a recurrence, and it is how a binomial coefficient is actually
computed.} \code{math.comb} does not form $n!$ and then divide it away, which
would build a number with hundreds of digits to answer a question whose answer
has four; it works down a triangle of additions, and every intermediate value
is a count of something.

The book's own habit applies to the derivation as well as to the arithmetic:
the argument above is a \emph{yes-or-no decision on one element}, which is
exactly the move Program~\ref{prog:F10} used to get $2^{n}$. The same move
answers a second question here, which is the next frame.
\end{fr}

\begin{fr}
Add up a whole row. What is
$\binom{\val{p12.abl.n}}{0} + \binom{\val{p12.abl.n}}{1} + \cdots +
\binom{\val{p12.abl.n}}{\val{p12.abl.n}}$, and why?
\dotline
\nextframe
\end{fr}

\begin{fr}
\ans{$\val{p12.abl.all}$}
Which is $2^{\val{p12.abl.n}}$, and is Program~\ref{prog:F10}'s subset count
to the digit.

The reason is that the row \emph{is} the subsets, sorted by size. To choose a
subset you may pick its size and then pick which elements \dash{} that is the
left-hand side \dash{} or you may go through the elements deciding in or out,
which is the right-hand side. One collection of objects counted two ways, and
the identity falls out with no algebra at all.

\begin{note}
This is the second of the three gates against Program~\ref{prog:F10}. The
script sums that row and fails the build if the total is not the subset count
that program committed, so \enquote{choices} here and \enquote{subsets} there
cannot come apart.
\end{note}
\end{fr}

\begin{fr}
The second rule that is not a product is one you have in its two-set form.

Program~\ref{prog:F10} had two evaluation sets of $\val{p12.ie.a}$ and
$\val{p12.ie.b}$ examples sharing $\val{p12.ie.ab}$, and made the point that
adding the sizes counts the shared part twice.

Those are the same two sets. A third has arrived, of $\val{p12.ie.c}$
examples: it shares $\val{p12.ie.ac}$ with the first and $\val{p12.ie.bc}$
with the second, and $\val{p12.ie.abc}$ examples are in all three. How many
distinct examples is that?
\dotline
\nextframe
\end{fr}

\begin{fr}
\ans{$\val{p12.ie.union}$}
Not $\val{p12.ie.naive}$, which is $\val{p12.ie.over}$ too many.

\[ \lvert A \cup B \cup C \rvert = \val{p12.ie.naive}
   - \val{p12.ie.pairsum} + \val{p12.ie.abc} = \val{p12.ie.union} \]

The first two terms are the rule you had. The third is the new part and it is
the one worth understanding rather than memorising: an example in all three
sets is counted \emph{three} times by the first term, and then subtracted
\emph{three} times by the second, once in each pairwise overlap. So it has
been removed altogether, and it has to come back once.

\textbf{The signs alternate because each round of correction overshoots the
last one.} That is the general rule, and it is called inclusion--exclusion for
exactly that reason.

\begin{note}
The script builds these as three real sets out of their seven regions and
compares the formula against the actual union, rather than checking the
arithmetic. It also gates the first two sets against
Program~\ref{prog:F10}'s committed sizes, so this is the same worked example
continued and not a new one that happens to use round numbers.
\end{note}
\end{fr}

\begin{fr}
For $n$ sets the rule has one term per non-empty subset of them, so
$2^{n} - 1$ terms: $\val{p12.ie.terms3}$ for three sets, and
$\val{p12.ie.terms10}$ for $\val{p12.ie.sets}$.

\begin{trapbox}
\textbf{Nobody uses inclusion--exclusion past three sets, and the reason is in
that number.} At $\val{p12.ie.sets}$ data sources it wants
$\val{p12.ie.terms10}$ intersections, each of which is itself a computation,
and every one of them must be right.

What people do instead is build the union and take its size, which is one
pass over the data and no correction terms at all \dash{} and that is the
right instinct, because $2^{n}-1$ terms is $2^{n}$ subsets minus the empty
one, so the rule costs exactly as much as enumerating what it was meant to
avoid enumerating.

Two and three sets are the cases that come up in a design review, and both are
one line. Keep those; reach for a set library for the rest.
\end{trapbox}
\end{fr}

\begin{fr}
The third rule counts nothing at all. It says what \emph{must} be true, which
is a different kind of statement and turns out to be the more useful one.

You hash $\val{p12.pig.items}$ documents into $\val{p12.pig.buckets}$ buckets.
Say something that is certainly true about the result, whatever the hash
function is and however the documents fall.
\dotline
\nextframe
\end{fr}

\begin{fr}
\ans{Some bucket holds at least $\val{p12.pig.floor}$ documents}
Because if every bucket held at most $\val{p12.pig.floor} - 1$, the total
would be at most $\val{p12.pig.buckets} \times 3 = 768$, and there are
$\val{p12.pig.items}$.

That is the \textbf{pigeonhole principle}: with $m$ items in $n$ boxes some
box holds at least $\lceil m/n \rceil$. The proof is the one just given, and
it is worth noticing that it is a proof by contradiction with no arithmetic in
it beyond a multiplication.

\begin{note}
The script tests this over $\val{p12.pig.trials}$ random assignments, and
\emph{tests} is the right word. A random sweep here cannot discover anything:
a counterexample would refute a proof, so what the sweep is checking is the
code, exactly as in Program~\ref{prog:P04}.
\end{note}

It also gives the first honest answer about hash collisions. With more
documents than buckets a collision is not likely, it is \textbf{certain} --
so a $\val{p12.bits64}$-bit hash cannot keep more than
$\val{p12.hash.n64}$ documents apart, whatever else is true.
\end{fr}

\begin{fr}
And that bound is almost useless, which is the point of the next section.

Nobody has $\val{p12.hash.n64}$ documents. The number that decides whether a
hash is big enough is not the one at which collisions become certain, it is
the very much smaller one at which they become \emph{likely} \dash{} and how
much smaller is the single most surprising number in this program.
\end{fr}

% ============================================================
\section{How many bits a hash needs}
\label{sec:P12-birthday}
\index{counting!birthday problem}
\index{hashing!collision}

\begin{fr}
Start with the version everybody has heard, because getting it wrong is what
makes the engineering version land.

$\val{p12.bday.m}$ people are in a room. What is the chance that two of them
share a birthday? Guess before working it out; there are
$\val{p12.bday.n}$ days.
\dotline
\nextframe
\end{fr}

\begin{fr}
\ans{$\val{p12.bday.p}$ per cent}
Rather more than half, from $\val{p12.bday.m}$ people and
$\val{p12.bday.n}$ days.

Most people guess something near $\val{p12.bday.naive}$ per cent, and that
guess is $\val{p12.bday.m}$ divided by $\val{p12.bday.n}$ \dash{} which is a
perfectly good answer to a different question: the chance that somebody shares
\emph{your} birthday.

\begin{trapbox}
\textbf{The question is not you against the room. It is every pair in the room
against every other.} $\val{p12.bday.m}$ people make
$\binom{\val{p12.bday.m}}{2} = \val{p12.bday.pairs}$ pairs, and each pair is
its own chance to match.

$\val{p12.bday.pairs}$ chances at one in $\val{p12.bday.n}$ apiece is
manifestly not a small number of chances, and that is the whole of the
surprise. The count of people grows one at a time; the count of pairs grows
with the square, which is Program~\ref{prog:F10}'s formula and its
\enquote{twice the data is four times the work} arriving in an entirely
different subject.

Everybody who is surprised by this is surprised for the same reason: they
counted the wrong objects. Count the pairs and there is nothing left to be
surprised by.
\end{trapbox}

\mermaidfig{p12-pairs-not-items}{Why the count that matters is the pairs, and
what that does to the size of hash you need.}{pairs, not items}
\end{fr}

\begin{fr}
The exact statement is a product. Nobody has matched so far, so the second
person must miss the first, the third must miss both, and so on:
\[ P(\text{no match}) = \prod_{i=0}^{m-1}\left(1 - \frac{i}{N}\right) \]
which is Program~\ref{prog:F04}'s pi, and is also a recurrence \dash{} each
term is the one before it times one more factor.

There is a second expression for the same thing, and it comes straight from
the pair count. Each of the $\binom{m}{2}$ pairs misses with probability
$1 - 1/N$, so
\[ P(\text{a match}) \approx 1 - e^{-m(m-1)/2N} \]
using Program~\ref{prog:F07}'s exponential to turn a product of near-ones into
a single term.

Two formulas for one quantity, one exact and one not. Which should the code
use?
\nextframe
\end{fr}

\begin{fr}
\begin{ansblock}
It depends on the regime, and neither is the safe default.
\end{ansblock}

Measured, as relative error against a reference that is neither of them:

\begin{center}
\begin{tabular}{llll}
\toprule
 & \textbf{the product} & \textbf{the exponential} \\
\midrule
$\val{p12.bday.m}$ people, $\val{p12.bday.n}$ days
  & $\val{p12.reg.bday.naive}$ & $\val{p12.reg.bday.closed}$ \\
$\val{p12.check.m}$ hashes, $\val{p12.bits64}$ bits
  & $\val{p12.reg.mid.naive}$ & $\val{p12.reg.mid.closed}$ \\
$\val{p12.trap.m}$ hashes, $\val{p12.bits128}$ bits
  & $\val{p12.reg.tiny.naive}$ & $\val{p12.reg.tiny.closed}$ \\
\bottomrule
\end{tabular}
\end{center}

\textbf{Each one is right exactly where the other fails, and they fail for
opposite reasons.} The exact product is exact on the birthday case and returns
a relative error of \emph{one} at $\val{p12.bits128}$ bits, which is to say it
returns nothing at all. The exponential is $\val{p12.reg.bday.closed}$ out on
the birthday case \dash{} it gives $\val{p12.bday.closed}$ per cent where the
answer is $\val{p12.bday.p}$ \dash{} and is right to fifteen figures at hash
scale.
\end{fr}

\begin{fr}
The exponential's error is easy to account for: it treats the pairs as
independent, and they are not quite, and that matters only when there are few
enough of them to notice.

The product's failure is Program~\ref{prog:P02} arriving in a place nobody
expects it. Here is the whole of it:

\transcript{p12-collision-zero}

Every factor $1 - i/N$ is a number within one part in $10^{33}$ of one, and
Program~\ref{prog:P01} says what float64 does with those: it rounds each to
exactly $1$. The product is therefore exactly $1$, and subtracting it from $1$
has nothing left to say.

\begin{trapbox}
\textbf{A collision probability of exactly zero reads as a proof of safety,
and it is a rounding artefact.} There is no exception, no warning and no
\code{nan} \dash{} the function returns a float, the float is \code{0.0}, and
\code{0.0} is a completely reasonable-looking answer to the question that was
asked.

Program~\ref{prog:P02} named this and gave the rule: never form a probability
by subtracting a near-one from one. The version of it here is worse than the
one that program measured, because there the loss was digits and here it is
all of them, and because the wrong answer is the one that ends the
conversation.

Note the middle row of the table as well. At $\val{p12.bits64}$ bits and only
$\val{p12.check.m}$ items the product has \emph{already} lost five significant
figures, long before it fails outright. A formula does not stop working at a
threshold; it degrades, and the degradation is invisible until somebody
computes the same thing another way.
\end{trapbox}
\end{fr}

\begin{fr}
Now the engineering question. You are deduplicating a corpus by hashing:
$\val{p12.hash.m}$ documents, a $\val{p12.bits64}$-bit hash, so
$\val{p12.hash.n64}$ possible values.

What is the chance that two distinct documents get the same hash and one is
silently dropped?
\dotline
\nextframe
\end{fr}

\begin{fr}
\ans{About $\val{p12.hash.p64}$ per cent}
Rather than the $\val{p12.hash.naive64}$ that the number of documents over the
number of hash values suggests. \textbf{The two differ by a factor of
$\val{p12.hash.ratio}$}, and the difference is entirely the pairs.

\begin{aibox}
$\val{p12.hash.p64}$ per cent is a chance you would not accept anywhere else in
a pipeline. It is one corpus in forty containing at least one pair of distinct
documents that the deduplication step declared identical \dash{} and the one
it keeps is arbitrary, so the loss is a document, silently, with no log line
and nothing downstream that could notice.

The reason $\val{p12.bits64}$ feels ample is the calculation nobody does
aloud: $\val{p12.hash.n64}$ is an enormous number and
$\val{p12.hash.m}$ documents is a small fraction of it. Both statements are
true and neither is the question.
\end{aibox}
\end{fr}

\begin{fr}
So how large does the corpus have to be before a $\val{p12.bits64}$-bit hash
is a coin flip?

The pair count is about $m^{2}/2$, so a match becomes likely once
$m^{2}/2N$ is about $1$, which puts $m$ at about $\sqrt{N}$. Doing it properly
gives $m = \val{p12.half.c}\sqrt{N}$ for exactly a half.

\begin{center}
\begin{tabular}{lll}
\toprule
\textbf{hash} & \textbf{values} & \textbf{items for a coin flip} \\
\midrule
$\val{p12.bits64}$-bit & $\val{p12.hash.n64}$ & $\val{p12.hash.half64}$ \\
$\val{p12.bits128}$-bit & $\num{3.40e38}$ & $\val{p12.hash.half128}$ \\
\bottomrule
\end{tabular}
\end{center}

\textbf{A $\val{p12.bits64}$-bit hash is a coin flip at
$\val{p12.hash.half64}$ documents}, which is a corpus that exists. At
$\val{p12.bits128}$ bits the same threshold is $\val{p12.hash.half128}$, and
at $\val{p12.hash.m2}$ documents the collision probability is
$\val{p12.hash.p128}$.

That is the whole argument for $\val{p12.bits128}$ bits, and it is not that
$\val{p12.bits128}$ is twice $\val{p12.bits64}$. \textbf{Doubling the bits
squares the number of values and therefore squares the corpus you can hold},
because the safe capacity goes with the square root: $2^{32}$ against
$2^{64}$.
\end{fr}

\begin{fr}
The rule to carry out of this section is one line, and it is worth more than
the formula it came from.

\textbf{A hash's safe capacity is about the square root of the number of
values it can take}, because what has to stay distinct is the pairs and there
are about $m^{2}/2$ of those. Halve the exponent and you have the answer:
$\val{p12.bits64}$ bits holds about $2^{32}$ items, $\val{p12.bits128}$ bits
about $2^{64}$.

And the two statements this section has made about collisions are worth
keeping apart, because they answer different questions. The pigeonhole
principle says collisions become \textbf{certain} at
$\val{p12.pig.certain64}$ items. The birthday calculation says they become
\textbf{likely} at $\val{p12.hash.half64}$. Those differ by a factor of about
four billion, and it is always the second one that decides the design.
\end{fr}

% ============================================================
\section{Two counts that decide a design}
\label{sec:P12-designs}
\index{counting!search space}
\index{counting!Shapley value}

\begin{fr}
Two more counts, and neither needs anything the last four sections have not
given. What they need is the willingness to compute them before proposing the
thing they price.

A decoder generates $\val{p12.beam.len}$ tokens from a vocabulary of
$\val{p12.beam.vocab}$. How many distinct sequences could it produce? You will
want Program~\ref{prog:F03} for this.
\dotline
\nextframe
\end{fr}

\begin{fr}
\ans{$\val{p12.beam.space}$}
Which is $\val{p12.beam.vocab}^{\val{p12.beam.len}}$, the top-left cell of
this program's table \dash{} order matters and tokens certainly repeat.

Taking $\logb{10}$ of it is the only way to see what it is:
$\val{p12.beam.len} \times \logb{10}\val{p12.beam.vocab}$, which is a shade
over ninety. A ninety-digit number, for twenty tokens.

Now a beam search of width $\val{p12.beam.a}$ runs over that space. At each of
the $\val{p12.beam.len}$ steps it extends each of its
$\val{p12.beam.a}$ live sequences by every one of the
$\val{p12.beam.vocab}$ tokens and keeps the best $\val{p12.beam.a}$. How many
continuations does it score, in total?
\nextframe
\end{fr}

\begin{fr}
\ans{$\val{p12.beam.scored.a}$}
From $\val{p12.beam.len} \times \val{p12.beam.a} \times \val{p12.beam.vocab}$.
As a fraction of the space, that is $\val{p12.beam.frac.a}$.

Double the beam to $\val{p12.beam.b}$ and the work doubles, to
$\val{p12.beam.scored.b}$ scored continuations. The fraction of the space seen
also doubles, to $\val{p12.beam.frac.b}$.

\begin{warning}
\textbf{What this count rules out is an argument, not a practice.} Beam search
demonstrably produces better output than greedy decoding, and this book has
not measured by how much or where it stops helping.

What the count settles is \emph{why} it cannot be coverage. Going from
$\val{p12.beam.a}$ to $\val{p12.beam.b}$ takes the fraction of the space
examined from $\val{p12.beam.frac.a}$ to $\val{p12.beam.frac.b}$, and both of
those are zero for every practical purpose. A wider beam cannot be helping by
seeing more of the space, because no width anybody would run sees any of it.

Whatever it buys, it buys from the \emph{ordering} the model imposes on the
space \dash{} the good sequences are not scattered through those ninety digits
at random, they are concentrated where the model's own scores are high, and a
wider beam hedges against the score being locally misleading. That is a claim
about the model rather than about the search, and it is the claim to argue
with when somebody proposes doubling the beam.
\end{warning}
\end{fr}

\begin{fr}
The last count is the one that should change how you read a plot.

A Shapley value attributes a model's output to its features, and the
definition is a counting statement: \textbf{a feature's attribution is its
average marginal contribution over every ordering of the features.} Add the
features one at a time in some order, note how much this one added when its
turn came, and average that over all the orders.

For $\val{p12.shap.big}$ features, how many orderings is that?
\dotline
\nextframe
\end{fr}

\begin{fr}
\ans{$\val{p12.shap.big.ord}$}
Which is $\val{p12.shap.big}!$ and settles nothing, because it is not
computable.

But look again at what the definition needs. When a feature's turn comes, what
it adds depends on the set of features already there \dash{} and not at all on
the order they were added in. So every ordering that puts the same set before
it gives the same marginal contribution, and they can be collected.

Collecting them replaces $\val{p12.shap.big}!$ orderings with
$2^{\val{p12.shap.big}} = \val{p12.shap.big.coal}$ subsets, each weighted by
how many orderings it stood for. \textbf{A factor of
$\val{p12.shap.big.ratio}$}, for free, and it is the only reason the value can
be computed at all.

\mermaidfig{p12-orderings-or-subsets}{How an average over $n!$ orderings
collapses into a sum over $2^{n}$ subsets.}{orderings collapse to subsets}
\end{fr}

\begin{fr}
\begin{note}
The identity is proved by enumeration in \code{code/p12\_combinatorics.py}
rather than asserted: with $\val{p12.shap.n}$ features it computes every
feature's value both ways \dash{} averaging over all
$\val{p12.shap.orderings}$ orderings, and summing over all
$\val{p12.shap.coalitions}$ subsets with the weights \dash{} and requires them
to be equal as exact fractions.

The same six-feature example makes the other point attribution is for. Two of
its features cover exactly the same ground. Each of them is worth
$\val{p12.shap.alone}$ on its own and each is attributed
$\val{p12.shap.redundant}$, while a feature covering ground nothing else
reaches is attributed the whole $\val{p12.shap.top}$ it is worth. The
attributions add up to the $\val{p12.shap.total}$ the full set achieves, which
is the property that makes them shares rather than scores.
\end{note}

So: $\val{p12.shap.big.coal}$ evaluations for $\val{p12.shap.big}$ features,
which a computer will do. Now $\val{p12.shap.huge}$ features. How many?
\nextframe
\end{fr}

\begin{fr}
\ans{$\val{p12.shap.huge.coal}$}
At a millisecond apiece \dash{} which is optimistic for anything with a model
in it \dash{} that is $\val{p12.shap.huge.years}$ years.

\begin{trapbox}
\textbf{Exact Shapley attribution is exponential by construction, so every
implementation you have used is a sampling approximation.} This is not a
shortcoming of any particular library. The quantity is defined as an average
over $2^{n}$ subsets, and there is no rearrangement that makes it fewer,
because each of those subsets can genuinely give a different answer.

What follows is worth carrying into the next attribution plot you are shown.
The bars have a sampling error, that error is not usually drawn, and it is
largest exactly where the features interact \dash{} which is where the plot is
being used to make an argument. Two runs of the same explainer on the same
input can reorder the middle of the ranking.

The question to ask of such a plot is the one this program has been asking all
along, and it is a counting question: \textbf{how many of the
$2^{n}$ were actually evaluated?}
\end{trapbox}
\end{fr}

\begin{fr}
One closing observation, which is what this program has really been about.

Every count in it was a decision made visible. Whether a hash is wide enough,
whether an ablation study fits in a week, whether a beam can be widened, how
far an attribution can be trusted \dash{} each was settled by a number that
took one multiplication and that almost nobody computes before committing to
the thing it prices.

\textbf{The rules are not the useful part.} There are four of them and they
fit on a card. What is useful is the habit of noticing that a plan contains a
count at all, and Program~\ref{prog:F10} put it best: write the count down
before writing the loop.
\end{fr}

% ============================================================
\begin{summarybox}
\sumitem{1--2}{\result{\textbf{Two questions decide which rule applies}: does the
  order matter, and may anything repeat? \enquote{Pick $\val{p12.four.k}$ of
  $\val{p12.four.n}$} sounds complete and has three defensible answers until
  both are settled}.}
\sumitem{3--4}{\result{Two of the four counts, of $\val{p12.four.k}$ from
  $\val{p12.four.n}$: $n^{k} = \val{p12.four.ordrep}$ ordered with repeats,
  and $\frac{n!}{(n-k)!} = \val{p12.four.ordnorep}$ ordered without, where the
  pool shrinks by one at each step}.}
\sumitem{4--5}{\result{The product rule needs \textbf{the number of options at each
  step to be fixed, not the options themselves}, which is why a shrinking pool
  multiplies anyway. Program~\ref{prog:F10}'s stronger statement is the one
  for checking a grid; this weaker one is what lets an arrangement count be a
  product at all}.}
\sumitem{6}{\result{$\binom{n}{k} = \val{p12.four.unordnorep}$ is the arrangements
  divided by the $k! = \val{p12.kfact}$ orders each selection was counted in,
  and \textbf{that division is Program~\ref{prog:F10}'s \enquote{halve it}
  generalised} \dash{} it halved because $2! = 2$}.}
\sumitem{7--8}{\result{The fourth cell, unordered with repeats, is
  $\binom{n+k-1}{k} = \val{p12.four.unordrep}$, stated and not derived. The
  four counts of the same ten things taken three at a time span a factor of
  more than eight}.}
\sumitem{9--10}{\result{$\val{p12.four.n}!$ is $\val{p12.fact.n}$ orderings, and
  $n! = \frac{n!}{(n-k)!} \times (n-k)!$ is the shrinking pool read backwards
  \dash{} which is why an arrangement count is written as a quotient and
  computed as a product}.}
\sumitem{11--12}{\result{$\binom{n}{k} = \frac{n!}{k!(n-k)!}$, \emph{symbol Newtona} in
  Polish. $\binom{n}{2} = \frac{n(n-1)}{2}$ is Program~\ref{prog:F10}'s pair
  count, and the build fails if the two disagree}.}
\sumitem{13--14}{\result{\textbf{$\binom{n}{k} = \binom{n}{n-k}$, because choosing what
  to keep is choosing what to drop.} Both are $\val{p12.sym.val}$ at
  $\val{p12.sym.n}$ taken $\val{p12.sym.k}$ or $\val{p12.sym.rest}$ at a
  time}.}
\sumitem{15--16}{\result{An ablation over every subset of $\val{p12.abl.n}$ features
  is $\val{p12.abl.all}$ runs; capped at $\val{p12.abl.k}$ features it is
  $\val{p12.abl.upto}$, a factor of $\val{p12.abl.frac}$}.}
\sumitem{17--20}{\result{\textbf{Pascal's rule
  $\binom{n}{k} = \binom{n-1}{k-1} + \binom{n-1}{k}$} is a recurrence and is
  how the coefficient is computed; summing a row gives
  $\val{p12.abl.all} = 2^{\val{p12.abl.n}}$, because choosing a size and then
  which is the same act as deciding in-or-out per element}.}
\sumitem{21--23}{\result{\textbf{Inclusion--exclusion: the signs alternate because each
  round of correction overshoots the last.} Three sets of
  $\val{p12.ie.a}$, $\val{p12.ie.b}$ and $\val{p12.ie.c}$ hold
  $\val{p12.ie.union}$ distinct examples, not $\val{p12.ie.naive}$. It has
  $2^{n}-1$ terms, so past three sets build the union instead}.}
\sumitem{24--26}{\result{\textbf{Pigeonhole: $m$ items in $n$ boxes put at least
  $\lceil m/n \rceil$ somewhere}, whatever the hash does. It makes collisions
  \emph{certain} at $\val{p12.pig.certain64}$ items, which is not the number
  that decides anything}.}
\sumitem{27--28}{\result{\textbf{The birthday count is over pairs, not items.}
  $\val{p12.bday.m}$ people make $\val{p12.bday.pairs}$ pairs and share a
  birthday with probability $\val{p12.bday.p}$ per cent, against the
  $\val{p12.bday.naive}$ per cent that answers a different question}.}
\sumitem{29--31}{\result{Two formulas, each right where the other fails: the exact
  product is exact on the birthday case and returns \textbf{exactly zero} at
  $\val{p12.bits128}$ bits, while the exponential is
  $\val{p12.reg.bday.closed}$ out there and right to fifteen figures at hash
  scale. Program~\ref{prog:P02}'s rule, in a hash}.}
\sumitem{32--35}{\result{\textbf{A hash's safe capacity is about the square root of its
  number of values.} $\val{p12.hash.m}$ documents in $\val{p12.bits64}$ bits
  collide with probability $\val{p12.hash.p64}$ per cent \dash{}
  $\val{p12.hash.ratio}$ times the figure the naive reading gives \dash{} and
  the coin flip is at $\val{p12.hash.half64}$ documents against
  $\val{p12.hash.half128}$ for $\val{p12.bits128}$ bits}.}
\sumitem{36--38}{\result{$\val{p12.beam.vocab}^{\val{p12.beam.len}}$ is
  $\val{p12.beam.space}$ sequences and a beam of $\val{p12.beam.a}$ scores
  $\val{p12.beam.scored.a}$ of them. \textbf{Doubling the beam doubles a
  fraction of $\val{p12.beam.frac.a}$}, so coverage cannot be why a wider beam
  helps}.}
\sumitem{39--43}{\result{\textbf{A Shapley value averages over every ordering, and the
  orderings collapse to $2^{n}$ subsets} because only which features came
  before matters. $\val{p12.shap.big.coal}$ at $\val{p12.shap.big}$ features
  and $\val{p12.shap.huge.coal}$ at $\val{p12.shap.huge}$, so every
  implementation samples \dash{} and the bars carry an error nobody plots}.}
\end{summarybox}

\canyou

\begin{testexercises}
\item From a vocabulary of $50$ tokens, how many sequences of $4$ tokens are
  there? And how many \emph{sets} of $4$ distinct tokens?
  \answerto{$50^{4} = \num{6250000}$ sequences, since order matters and tokens
  repeat. And $\binom{50}{4} = \num{230300}$ sets, dividing the
  $50 \times 49 \times 48 \times 47$ arrangements by $4! = 24$. A factor of
  more than $27$ between them, decided by two words.}
\item Two data sources hold $600$ and $900$ records and share $150$; a third
  holds $400$, sharing $80$ with the first, $120$ with the second, and $40$
  with both. How many distinct records?
  \answerto{$1900 - 350 + 40 = 1590$. Sum $600 + 900 + 400 = 1900$, subtract
  $150 + 80 + 120 = 350$, add the triple overlap back once because subtracting
  the three pairs removed it three times after it had been counted three
  times.}
\item You hash $10$ million documents into a $32$-bit space. Give the
  approximate collision probability, and say what makes it that large.
  \answerto{Essentially $1$: the pair count is about $5 \times 10^{13}$
  against $\num{4.29e9}$ hash values, so the exponent
  $m(m-1)/2N$ is more than ten thousand. The coin flip for $32$ bits is at
  about $\num{7.7e4}$ documents \dash{} $\val{p12.half.c}\sqrt{2^{32}}$
  \dash{} so ten million is far past it.}
\item A grid search has $6$ models, $4$ learning rates and $3$ batch sizes,
  but two of the models cannot run at the largest batch size. How many runs?
  \answerto{$6 \times 4 \times 3 = 72$ less the forbidden ones. The two models
  lose one batch size each across all $4$ learning rates, so $8$ go:
  $64$ runs. Program~\ref{prog:F10}'s condition on the product rule, and the
  reason to count the exclusions when the grid is written down.}
\item Why is $\binom{n}{k}$ computed by Pascal's rule rather than from the
  factorials?
  \answerto{Because the factorials are enormous and cancel almost entirely:
  $\binom{50}{4}$ has six digits and $50!$ has sixty-five. Pascal's rule works
  down a triangle of additions in which every intermediate value is itself a
  count, so nothing is built that then has to be divided away.}
\end{testexercises}

\begin{furtherproblems}
\item Show that $\binom{n}{k} = \binom{n}{n-k}$ from the formula, in one line.
  \answerto{$\binom{n}{n-k} = \frac{n!}{(n-k)!\,(n-(n-k))!} =
  \frac{n!}{(n-k)!\,k!}$, which is $\binom{n}{k}$ with the two factors of the
  denominator written in the other order. The combinatorial argument is
  shorter still: choosing $k$ to keep is choosing $n-k$ to drop.}
\item A beam of width $b$ over $L$ steps and a vocabulary of $V$ scores
  $LbV$ continuations. What width would be needed to score one per cent of the
  space, at $V = \val{p12.beam.vocab}$ and $L = \val{p12.beam.len}$?
  \answerto{$b = \num{0.01}\,V^{L}/(LV)$, which is about $2 \times 10^{84}$. The
  point of the exercise is that the number is not merely large but
  unreachable, so no width is a coverage argument.}
\item The empty product gives $0! = 1$. Show that this is forced by Pascal's
  rule rather than being a convention chosen for convenience.
  \answerto{$\binom{n}{n} = \binom{n-1}{n-1} + \binom{n-1}{n}$, and the last
  term is zero because you cannot choose $n$ things from $n-1$. So
  $\binom{n}{n} = \binom{n-1}{n-1} = \cdots = \binom{0}{0} = 1$, and
  $\binom{n}{n} = \frac{n!}{n!\,0!}$ equals $1$ only if $0! = 1$. The same
  answer Program~\ref{prog:F04} reached from the accumulator a loop starts
  at.}
\item Why does doubling a hash from $\val{p12.bits64}$ to
  $\val{p12.bits128}$ bits square the corpus it can hold rather than doubling
  it?
  \answerto{Because the safe capacity goes with $\sqrt{N}$ and doubling the
  bits squares $N$: $\sqrt{N^{2}} = N$. So $2^{32}$ becomes $2^{64}$, and the
  exponent doubles because the square root halves the doubled exponent.}
\item A colleague proposes computing exact Shapley values for a model with
  $\val{p12.shap.huge}$ features by sampling orderings rather than subsets.
  Does that help?
  \answerto{Sampling is what every implementation does and it is the only
  option, but it makes the result an estimate rather than exact \dash{} so the
  word \emph{exact} has to go. Sampling orderings and sampling subsets are two
  routes to the same estimator, since the orderings sharing a set before a
  feature all give one contribution; the choice affects the variance and not
  what is being estimated.}
\end{furtherproblems}
